\section{Score-Based Models 与 DDPM 的统一视角}

\subsection{两条研究路径的汇合}

历史上，DDPM 和 score-based models 是两条独立发展的研究路径：

\begin{knowledgebox}{两条路径的发展历程}
\begin{itemize}
    \item \textbf{DDPM 路径}：Ho et al. (2020) 基于变分推断和 ELBO 推导出训练目标。从概率图模型和变分自编码器的传统出发，定义前向/逆向马尔可夫链，通过最大化 ELBO 来训练逆向过程。
    \item \textbf{Score-Based 路径}：Song \& Ermon (2019, 2020) 基于 score matching 和 Langevin dynamics。从统计学中的 score function 出发，通过 denoising score matching 训练 score estimator，用退火 Langevin dynamics 采样。
\end{itemize}
两条路径最终在 Song et al. (2021) 的 Score-Based SDE 框架中被统一。
\end{knowledgebox}

\subsection{等价关系的数学总结}

我们已经在前面的各节中逐步建立了所有等价关系，这里做一个系统性的总结：

\begin{enumerate}
    \item \textbf{训练目标等价}：DDPM 的 $\epsilon$-prediction 损失 $\equiv$ Denoising Score Matching 损失（差缩放因子）。
    \item \textbf{网络输出等价}：$\epsilon_\theta(x_t, t) = -\sqrt{1-\bar{\alpha}_t}\, s_\theta(x_t, t)$。
    \item \textbf{采样过程等价}：DDPM 的逆向采样 $\equiv$ 退火 Langevin dynamics（离散化形式）。
    \item \textbf{Tweedie 连接}：$\mathbb{E}[x_0 \mid x_t]$ 可以用 score function 表示，也可以用噪声预测器表示。
\end{enumerate}

\begin{importantbox}{统一视角的实际意义}
理解这种等价性不仅是学术上的优雅，更有重要的实际意义：
\begin{itemize}
    \item 可以借用 score-based 理论来分析 DDPM 的性质（如收敛性、采样质量）。
    \item 可以借用 DDPM 的实践技巧来改进 score-based 采样（如 schedule 设计）。
    \item Score function 的视角为条件生成（classifier guidance, classifier-free guidance）提供了自然的理论基础。
    \item Song et al. 的连续时间 SDE 框架将离散步骤推广到连续过程，开启了更灵活的采样器设计（如 DDIM, DPM-Solver 等）。
\end{itemize}
\end{importantbox}

\subsection{本章小结}

DDPM 和 score-based models 虽然从不同的出发点和动机发展而来，但在数学上完全等价。DDPM 提供了概率图模型和变分推断的视角，而 score-based models 提供了 score function 和 Langevin dynamics 的物理直觉。理解两种视角对于深入掌握扩散模型理论至关重要。

